Der Astronaut drückt mit derselben Kraft auf den Boden, mit welcher der Boden ihn stützt (Normalkraft, Wechselwirkungsgesetz). Steht er ruhig, ist die Normalkraft gleich der Gewichtskraft. Mit der Fallbeschleunigung auf dem Mond $g_\mathrm{M} = \gMoO$ (FoTaBe S.~210) ist
\begin{align}
 F_\mathrm{N} &= \FNoF = \mAs\cdot\gMo = \FNo \approx \result{\FNoP{0}{2}}
\end{align}

Kurzzeitig grösser wird die Kraft, wenn er seinen Körper nach oben beschleunigt, zum Beispiel beim Abspringen aus der Hocke (oder wenn er nach einem Sprung bei der Landung abbremst). Dann muss der Boden zusätzlich die beschleunigende Kraft liefern:
\begin{align}
 F_\mathrm{N} - m\,g_\mathrm{M} &= m\,a \quad\Rightarrow\quad F_\mathrm{N} = m\left(g_\mathrm{M}+a\right) > m\,g_\mathrm{M}
\end{align}
